\documentclass[11pt,a4paper]{letter}
\usepackage[margin=1in]{geometry}
\usepackage{parskip}
\signature{Yuanzhi He}
\address{School of Computational and Mathematical Sciences\\Cardiff University\\
Cardiff, United Kingdom\\heyuanzhi0728@163.com}
\begin{document}
\begin{letter}{The Editors\\\textit{Neurocomputing}}
\opening{Dear Editors,}

I submit for your consideration the manuscript \textit{Reference Matching for
Noise-Space Policy Fine-Tuning}.

Reinforcement learning over the latent noise input of a frozen generative policy
is an increasingly common way to raise the reward of an imitation-learned robot
policy, and the methods that try to preserve its behaviour almost all maximise
the entropy of a distribution over behavioural modes. The manuscript shows that
this rests on an assumption no one states: under a uniform reference, matching
the pretrained distribution and maximising entropy are the same objective, since
the Kullback-Leibler divergence to a uniform reference is $\ln K$ minus the
entropy. Behaviour cloning does not give uniform policies, so the two objectives
come apart in exactly the regime where these methods are used, and the
difference has gone unmeasured.

The paper replaces the entropy term with a total-variation penalty toward the
frozen policy's own mode histogram. It carries no learned component and one
scalar weight. Across 38 independently trained frozen policies and 1{,}235
fine-tuning runs on two manipulation tasks, with each frozen policy as the unit
of analysis, it ends closer to the pretrained behaviour than entropy
maximisation on both tasks, and the margin widens as the pretrained
distribution departs from uniform, which is the trend the identity predicts.

The work is original, is not under consideration elsewhere, and has no competing
interests to declare. The code, the run records behind every table and figure,
and the analysis plans registered before the runs are public.

\closing{Yours sincerely,}
\end{letter}
\end{document}
